\section*{Chapitre \ref{ch:pl:accelerators-for-quantitative-work} --- Accélérateurs pour le travail quantitatif}

\begin{solution}{exo:pl:accelerators-for-quantitative-work:1}
$2n^3$ opérations sur $3 \times 8n^2$ octets~: $n/12$. Le noyau est \omterm{def:pl:accelerators-for-quantitative-work:bound}{limité par le calcul} (compute-bound) lorsque $n/12$ dépasse le point de rupture~: à partir de $n = 34$ sur le cœur (point de rupture 2.79),
à partir de $n = 122$ sur le périphérique (point de rupture 10.1).
\end{solution}

\begin{solution}{exo:pl:accelerators-for-quantitative-work:2}
Son \omterm{def:pl:accelerators-for-quantitative-work:intensity}{intensité arithmétique} (inférieure à 1.3 opération par octet pour chaque noyau) est inférieure aux deux points de rupture, donc sur les deux machines son temps est égal aux octets divisés par la bande passante,
et le rapport des temps est le rapport des bandes passantes, 140~; le pic de performance, bien supérieur, du périphérique n'est jamais atteint.
\end{solution}

\begin{solution}{exo:pl:accelerators-for-quantitative-work:3}
$1/(0.05 + 0.95/140) = 17.6$~; avec 99~\%, $1/(0.01 + 0.99/140) = 58.6$. La part qui reste sur l'hôte fixe le plafond~: 20 et 100 pour
une accélération infinie.
\end{solution}

\begin{solution}{exo:pl:accelerators-for-quantitative-work:4}
Le \omterm{def:pl:accelerators-for-quantitative-work:roofline}{modèle roofline} compte \texttt{exp} comme une seule opération, mais une exponentielle prend de nombreux cycles~; le noyau est limité par le débit de
l'exponentielle, non par la bande passante. Une exponentielle vectorisée (une bibliothèque de fonctions transcendantes SIMD), le calcul de l'exponentielle de la
somme cumulée une seule fois par trajectoire plutôt qu'à chaque pas, ou la fusion du pas avec le noyau suivant pour que $S$ ne soit
pas écrit puis relu à nouveau, déplaceraient le noyau sur le graphique.
\end{solution}

\begin{solution}{exo:pl:accelerators-for-quantitative-work:5}
Le kilooctet de chaque transaction traverse la liaison à 128~Go/s (7.8~ns) et ses 8~ko de trafic mémoire coûtent 2.4~ns au périphérique~; le CPU a besoin de 334~ns par
transaction. Le rapport entre le temps du CPU par transaction et celui du périphérique vaut 32.7, quel que soit le lot, une fois la latence de lancement amortie.
\end{solution}

\begin{solution}{exo:pl:accelerators-for-quantitative-work:6}
La génération des trajectoires et le payoff à chaque date exécutent les mêmes instructions sur chaque trajectoire~: aucune \omterm{def:pl:accelerators-for-quantitative-work:divergence}{divergence des threads}. La décision d'exercice (poursuivre ou
exercer, selon la valeur de la trajectoire face à la régression) diverge d'une trajectoire à l'autre, et les trajectoires déjà exercées cessent de contribuer~; la
régression elle-même est un ensemble de réductions, sans problème sur le périphérique mais avec une synchronisation à chaque date.
\end{solution}

\begin{solution}{exo:pl:accelerators-for-quantitative-work:7}
Un mégaoctet par requête ajoute 7.8~microsecondes de transfert aux 10 de latence~: le seuil de rentabilité passe de 31 à 56 transactions par requête. Garder les
données de marché sur le périphérique et n'y envoyer que ce qui a changé vaut autant que diviser par deux la latence de lancement.
\end{solution}

\begin{solution}{exo:pl:accelerators-for-quantitative-work:8}
Le rapport des pics ne s'applique qu'aux calculs \omterm{def:pl:accelerators-for-quantitative-work:bound}{limités par le calcul} (compute-bound)~; un calcul de risque de Monte-Carlo et d'agrégation est \omterm{def:pl:accelerators-for-quantitative-work:bound}{limité par la mémoire} (memory-bound) et ne gagne au plus que le rapport des
bandes passantes, 140 (quatre minutes et non une), avant les transferts, les lancements et la part restée sur l'hôte (Amdahl). Et la comparaison se fait avec un seul cœur~: face à
la totalité du socket du serveur, le gain est plusieurs fois plus faible.
\end{solution}

\begin{solution}{pb:pl:accelerators-for-quantitative-work:1}
\begin{enumerate}
\item 66.9~GFLOP/s, 24.0~Go/s, un point de rupture de 2.79 opérations par octet.
\item 34~TFLOPS en double précision, 3.35~To/s, un point de rupture de 10.1~; 508 et 140 fois le cœur.
\item Le pas de trajectoire, le payoff du panier et l'agrégation des expositions sur les deux~; le produit matriciel est \omterm{def:pl:accelerators-for-quantitative-work:bound}{limité par le calcul} (compute-bound) sur les deux.
\item Il est limité par l'exponentielle, que la convention de comptage traite comme une seule opération.
\item La latence, les transferts, la \omterm{def:pl:accelerators-for-quantitative-work:divergence}{divergence des threads}, la précision, le coût du portage, le prix et la puissance.
\item 66 milliards d'opérations et 98~Go de trafic mémoire~; seules les descriptions des transactions traversent la liaison et le profil revient.
\item 4.09~secondes sur le cœur, 29.3~millisecondes sur le périphérique.
\item Il est \omterm{def:pl:accelerators-for-quantitative-work:bound}{limité par la mémoire} (memory-bound)~: 140 est le rapport des bandes passantes.
\item Au plus 17.6 fois si 95~\% du processus est déporté.
\item Le rapport par cœur se réduit au rapport entre la bande passante du périphérique et celle du socket.
\item 0.33~microseconde sur le cœur, 10.0 sur le périphérique~: 30 fois plus lent.
\item Au seuil de 31 transactions par requête.
\item 4, 16, 31 et 155 transactions pour des latences de lancement de 1, 5, 10 et 50~microsecondes.
\item La liaison~: au plus 32.7 fois.
\item Regrouper ses requêtes par lots (des centaines par appel) et garder les données de marché sur le périphérique.
\item \textbf{Résultat nommé.} Le Monte-Carlo d'exposition tourne 140 fois plus vite de bout en bout sur le périphérique que sur un seul cœur (17.6 fois pour le
processus entier, avec 95~\% déportés)~; une requête de prix pour une seule transaction tourne 30 fois plus lentement~; le périphérique atteint le seuil de rentabilité à 31 transactions par requête avec une
latence de lancement de 10~microsecondes.
\item Les calculs sur tableaux volumineux, par lots, \omterm{def:pl:accelerators-for-quantitative-work:bound}{limités par la mémoire} ou par le calcul, avec peu de données qui traversent la liaison~: le Monte-Carlo d'exposition et des XVA, la réévaluation de scénarios,
l'entraînement des modèles.
\item Sa bande passante soutenue et son débit en double précision sur nos noyaux, la latence de lancement et de transfert de nos requêtes, et l'effort de
portage d'un noyau réel.
\item Parce que les résultats de risque et de valorisation sont comparés et réconciliés en double précision~; la simple précision modifierait le plafond du périphérique et
devrait d'abord être validée.
\item Quand le calcul est volumineux, par lots, et reste assez longtemps sur le périphérique pour amortir le coût du transfert vers lui.
\end{enumerate}
\end{solution}

\begin{solution}{iq:pl:accelerators-for-quantitative-work:1}
Opérations par octet déplacé. En deçà du point de rupture d'une machine, un noyau est limité par la bande passante, au-delà par l'arithmétique~; un matériel avec un pic plus élevé
ne change rien pour un \omterm{def:pl:accelerators-for-quantitative-work:bound}{noyau limité par la mémoire}, et le rapport des bandes passantes est le gain maximal qu'il peut tirer d'une mémoire plus rapide.
\iqlookfor{Raisonnement roofline.}
\end{solution}

\begin{solution}{iq:pl:accelerators-for-quantitative-work:2}
Le Monte-Carlo est volumineux et amortit les coûts fixes du périphérique~; chaque requête de valorisation est petite, donc son temps est dominé par le lancement,
la synchronisation et le transfert, que le CPU ne paie pas. Regrouper les requêtes par lots ou les garder sur le CPU.
\iqlookfor{Coûts fixes face à la taille du calcul.}
\end{solution}

\begin{solution}{iq:pl:accelerators-for-quantitative-work:3}
Compter les opérations du noyau, les octets de mémoire et les octets transférés~; prendre le pic, la bande passante et la liaison du périphérique dans sa fiche technique, et retenir une
latence de lancement mesurée ou une hypothèse prudente~; borner son temps avec le roofline et les coûts fixes~; comparer au plafond
mesuré du CPU~; appliquer la loi d'Amdahl à l'ensemble du processus.
\iqlookfor{Un modèle de coûts avec hypothèses explicites.}
\end{solution}

\begin{solution}{iq:pl:accelerators-for-quantitative-work:4}
Avec un générateur à compteur, les nombres aléatoires ne dépendent que de la trajectoire, du pas et de la graine, de sorte que les deux machines tirent les mêmes~; il
faut alors que l'arithmétique soit menée avec la même précision et dans le même ordre (les réductions sont la différence habituelle), et que les fonctions transcendantes soient comparées à une
tolérance ou prises dans la même bibliothèque.
\iqlookfor{Générateurs à compteur~; réductions et précision.}
\end{solution}

\begin{solution}{iq:pl:accelerators-for-quantitative-work:5}
Inventorier les charges de travail avec leurs opérations, leurs octets, leurs transferts et leurs tailles de lots~; les placer sur des roofline mesurés et cités~; prédire les
gains de bout en bout avec les coûts fixes et la loi d'Amdahl~; porter un noyau de la charge la plus prometteuse et le mesurer~; décider à partir du gain mesuré,
en le comparant au portage, à la maintenance, au prix et à la consommation électrique, et garder une implémentation CPU comme référence.
\iqlookfor{Modéliser d'abord, mesurer un noyau, garder la référence.}
\end{solution}
